\section{Learned scalar fields with population errors}
\label{app:scalar-population}
This section gives the training and evaluation methods for the scalar
fields of Section~\ref{sec:scalar-learned}. We estimate population
errors by numerical integration over a Gaussian start.\AIassistance

\subsection{Target, fields, and training}
The target is $Y=\mu\xi+\sigma Z$ with $\mu=1.5$, $\sigma=0.5$, a
symmetric random sign $\xi$, and an independent standard Gaussian $Z$.
The path is the independent interpolation \eqref{eq:independent}. Thus
$p_t$ is an equal mixture of the two Gaussians
$\mathcal N(\pm\mu t,(1-t)^2+\sigma^2t^2)$. The dimension is odd, so
Theorem~\ref{thm:odd-transfer} applies. The exact conditional velocity is
\eqref{eq:mixture-generator} with $d=1$.

Each field is a network with the input $(x,t)$, two hidden layers of 48
tanh units, and a scalar output. It has 2,545 parameters. The step size
is not an input. There are three objectives.
\begin{itemize}
\item \emph{CFM.} The conditional flow matching loss
$\E|u(t,(1-t)G+tY)-(Y-G)|^2$, with $t$ uniform on $[0,1]$ and a
minibatch of 1,024 independent draws.
\item \emph{One-grid fit.} The CFM loss plus $10\,L_8$.
\item \emph{Two-grid fit.} The CFM loss plus $10\,(L_8+L_{16})/2$.
\end{itemize}
The loss $L_n$ is computed as follows. Run $n$ uniform Euler steps from
1,024 new Gaussian start points. Take each of the times
$1/4,1/2,3/4,1$. Sort the Euler outputs and sort 1,024 new samples of
$p_t$. Then take the mean squared difference of the sorted values. This value is the squared
$W_2$ distance between the two empirical distributions. The loss $L_n$
is the mean over the four times, and its gradient is taken through the
Euler steps.

Each run uses 8,000 Adam updates with the learning rate $10^{-3}$ in
double precision. New samples are drawn at each update. Each objective
uses the ten seeds 0 to 9, which gives \ScalarModelCount{} fields. All
settings were fixed before the evaluation, and no field was selected
with the evaluation results.

\subsection{Population errors by quadrature}
Let $z\mapsto\widehat X_k(z)$ be the Euler map from the start point to
the state at node $k$. We evaluate it on 8,001 equally spaced points
$z\in[-9,9]$. If the map is increasing, the pairing of equal quantiles
is optimal, and
\[
 W_2^2(\Law(\widehat X_k),p_t)
 =\int\big(\widehat X_k(z)-Q_t(\Phi(z))\big)^2\varphi(z)\dd z .
\]
Here $Q_t$ is the quantile function of $p_t$, and $\Phi$ and $\varphi$
are the standard Gaussian distribution function and density. We compute
$Q_t$ by bisection on the distribution function of the mixture, and the
integral by the trapezoidal rule. The Gaussian mass outside $[-9,9]$ is
below $3\times10^{-19}$. For all \ScalarModelCount{} fields, the Euler maps increase on the
quadrature nodes. We also bound the derivative of each update,
$1+\partial_xu/N$, from the network weights. For $N=4,8,16$ we
propagate bounds through both layers on intervals covering the real
line. For $N\ge32$ we use the sum of the negative products of weights
along paths from the input to the output. Each tanh derivative lies
between zero and one. The smallest computed lower bound is above
$0.037$. These checks use double precision without directed rounding. For the exact conditional
velocity, this quadrature gives the skew-field columns of
Table~\ref{tab:population-example}, which contain only the error in the mixture
coordinate, to seven digits.

The reference for the exact flow of a field is the classical
fourth-order Runge--Kutta method with 1,024 steps. The \emph{model bias}
is the largest $W_2$ distance between this reference and the target path
at 64 equally spaced times. The \emph{integration error} at $N$ steps compares the
Euler run and the reference from the same start points. It is their
largest root mean squared distance at the four shared times. The
\emph{target error} at $N$ steps is the largest $W_2$ distance to the
target path at the four shared times. The quantities $E_c$ and $E_f$ of
\eqref{eq:errors} use all nodes of the coarse grid.

\subsection{Results}
Table~\ref{tab:scalar-models} gives the results for each field. The
means are in Table~\ref{tab:scalar} of the main article.

\emph{CFM fields.} The target error is smaller at 16, 64, and 256 steps
than at 8 steps for all ten seeds. The mean is \ScalarCfmEight{} at 8
steps and \ScalarCfmLarge{} at 256 steps. The model bias is between
\ScalarCfmBiasMin{} and \ScalarCfmBiasMax{}, with a mean of
\ScalarCfmBias{}. For \ScalarCfmArgminLast{} of the ten seeds, the
smallest target error is at 256 steps. For the other seeds it is at an
intermediate step count, and the difference from the value at 256 steps
is small.

\emph{One-grid fields.} The mean target error is \ScalarOneEight{} at 8
steps, \ScalarOneSixteen{} at 16 steps, \ScalarOneSixtyFour{} at 64
steps, and \ScalarOneLarge{} at 256 steps. The smallest target error is
at 8 steps for \ScalarOneArgminEight{} seeds and at 16 steps for
\ScalarOneArgminSixteen{} seeds. Compare each value with the value
at 8 steps. The target error is larger at 16 steps for
\ScalarOneIncreaseSixteen{} of the ten seeds, at 64 steps for
\ScalarOneIncreaseSixtyFour{} seeds, and at 256 steps for
\ScalarOneIncreaseLarge{} seeds. The exact two-sided
Wilcoxon signed-rank tests for these three paired comparisons give
$p=\ScalarOneWilcoxonSixteen$, $p=\ScalarOneWilcoxonSixtyFour$, and
$p=\ScalarOneWilcoxonLarge$. The mean model bias is \ScalarOneBias{}.
The mean integration error decreases from \ScalarOneOwnEight{} at 8
steps to \ScalarOneOwnLarge{} at 256 steps.

\emph{Two-grid fields.} The mean target error is \ScalarTwoEight{} at 8
steps, \ScalarTwoSixteen{} at 16 steps, \ScalarTwoSixtyFour{} at 64
steps, and \ScalarTwoLarge{} at 256 steps. The smallest target error is
at 16 steps for \ScalarTwoArgminSixteen{} of the ten seeds. Compared
with the value at 8 steps, the target error is larger at 64 steps for
\ScalarTwoIncreaseSixtyFour{} seeds and at 256 steps for
\ScalarTwoIncreaseLarge{} seeds ($p=\ScalarTwoWilcoxonSixtyFour$ and
$p=\ScalarTwoWilcoxonLarge$). Compared with the value at 16 steps, it
is larger at 64 steps for \ScalarTwoFromSixteenIncreaseSixtyFour{} seeds
and at 256 steps for \ScalarTwoFromSixteenIncreaseLarge{} seeds
($p=\ScalarTwoFromSixteenWilcoxonSixtyFour$ and
$p=\ScalarTwoFromSixteenWilcoxonLarge$). The same comparison with the
value at 16 steps gives \ScalarOneFromSixteenIncreaseSixtyFour{} and
\ScalarOneFromSixteenIncreaseLarge{} seeds for the one-grid fields, and
\ScalarCfmFromSixteenIncreaseSixtyFour{} and
\ScalarCfmFromSixteenIncreaseLarge{} seeds for the CFM fields. The mean
model bias of the two-grid fields is \ScalarTwoBias{}.

For a fixed field, the target error at the four evaluation times
approaches the bias over those times as the step count increases. The
theoretical two-grid error uses all shared nodes and approaches twice
the bias over the full time interval. The table gives a 64-time estimate
of bias and the two-grid error at $N=8$. Theorem~\ref{thm:characterization}
concerns the best error over a class of fields. A fit at small step counts
decreases the target error there. For the one-grid fields it also
increases the model bias. The weight of the added loss was fixed and
was not tuned, and a different weight can change the size of the
effect.

\begin{table}[ht]
\centering
\caption{Learned scalar fields: target error at $N$ Euler steps for each
field, model bias, and the two-grid error $E_c+E_f$ of \eqref{eq:errors}
at $N=8$. All values are population $W_2$ distances computed by
quadrature.}
\label{tab:scalar-models}
\begin{tabular}{@{}lrcccccc@{}}
\toprule
Field & Seed & $N=8$ & $16$ & $64$ & $256$ & Model bias & $E_c+E_f$ at $N=8$\\
\midrule
CFM & 0 & 0.146 & 0.072 & 0.025 & 0.021 & 0.023 & 0.218\\
CFM & 1 & 0.150 & 0.097 & 0.089 & 0.093 & 0.094 & 0.247\\
CFM & 2 & 0.143 & 0.078 & 0.050 & 0.051 & 0.052 & 0.220\\
CFM & 3 & 0.174 & 0.119 & 0.101 & 0.102 & 0.102 & 0.293\\
CFM & 4 & 0.148 & 0.076 & 0.040 & 0.040 & 0.041 & 0.225\\
CFM & 5 & 0.137 & 0.067 & 0.040 & 0.044 & 0.046 & 0.204\\
CFM & 6 & 0.257 & 0.217 & 0.198 & 0.196 & 0.195 & 0.473\\
CFM & 7 & 0.234 & 0.196 & 0.186 & 0.186 & 0.187 & 0.430\\
CFM & 8 & 0.187 & 0.143 & 0.131 & 0.132 & 0.133 & 0.330\\
CFM & 9 & 0.161 & 0.089 & 0.045 & 0.039 & 0.039 & 0.251\\
\addlinespace
One-grid fit & 0 & 0.073 & 0.079 & 0.110 & 0.119 & 0.123 & 0.152\\
One-grid fit & 1 & 0.068 & 0.068 & 0.099 & 0.108 & 0.112 & 0.136\\
One-grid fit & 2 & 0.088 & 0.090 & 0.114 & 0.122 & 0.125 & 0.178\\
One-grid fit & 3 & 0.107 & 0.112 & 0.138 & 0.147 & 0.150 & 0.219\\
One-grid fit & 4 & 0.122 & 0.129 & 0.151 & 0.159 & 0.161 & 0.251\\
One-grid fit & 5 & 0.168 & 0.169 & 0.185 & 0.191 & 0.193 & 0.336\\
One-grid fit & 6 & 0.065 & 0.041 & 0.068 & 0.078 & 0.084 & 0.106\\
One-grid fit & 7 & 0.055 & 0.047 & 0.082 & 0.092 & 0.096 & 0.102\\
One-grid fit & 8 & 0.053 & 0.049 & 0.084 & 0.094 & 0.098 & 0.102\\
One-grid fit & 9 & 0.048 & 0.057 & 0.096 & 0.107 & 0.110 & 0.105\\
\addlinespace
Two-grid fit & 0 & 0.091 & 0.090 & 0.114 & 0.122 & 0.125 & 0.181\\
Two-grid fit & 1 & 0.100 & 0.091 & 0.110 & 0.117 & 0.120 & 0.191\\
Two-grid fit & 2 & 0.112 & 0.103 & 0.119 & 0.125 & 0.127 & 0.215\\
Two-grid fit & 3 & 0.107 & 0.073 & 0.079 & 0.085 & 0.088 & 0.180\\
Two-grid fit & 4 & 0.103 & 0.092 & 0.109 & 0.116 & 0.118 & 0.196\\
Two-grid fit & 5 & 0.079 & 0.050 & 0.070 & 0.079 & 0.082 & 0.129\\
Two-grid fit & 6 & 0.086 & 0.064 & 0.081 & 0.088 & 0.091 & 0.150\\
Two-grid fit & 7 & 0.077 & 0.041 & 0.059 & 0.068 & 0.073 & 0.117\\
Two-grid fit & 8 & 0.089 & 0.048 & 0.053 & 0.059 & 0.062 & 0.137\\
Two-grid fit & 9 & 0.078 & 0.039 & 0.057 & 0.067 & 0.071 & 0.117\\
\bottomrule
\end{tabular}

\end{table}

\section{Variability of sample-based comparisons across evaluation sets}
\label{app:set-variability}
Section~\ref{app:experiments} reports the sliced distance SW2 on one
evaluation set of 1,024 samples and on one evaluation set of 262,144
samples. The two sets give different directions for the change between
8 and 16 steps. Here we measure how often the direction changes across
independent evaluation sets of the same size.\AIassistance

\subsection{Protocol}
The nine two-dimensional fields of Section~\ref{app:experiments} and its
128 projection directions are fixed. Only
the evaluation samples change. One evaluation set contains start points
for the numerical runs, and independent noise and data points for the
target path. We draw \VarSmallSets{} independent sets of 1,024 samples and
\VarMediumSets{} independent sets of 16,384 samples. Their random seeds are different from the seeds of
those two sets. For each set we compute SW2 at 8, 16, and 64 Euler steps
as the maximum over the five diagnostic times, in double precision.
This implementation gives the values of Table~\ref{tab:seed-grid-details} on
the separate 1,024-sample set with differences below $3\times10^{-8}$.

\subsection{Results}
Table~\ref{tab:set-variability} gives the results for each field and
each size. It gives the mean and the standard deviation across sets of
SW2, of its change from 8 to 16 steps, and of its change from 8 to 64
steps. It also gives the percentage
of sets in which SW2 increases.

\emph{Change from 8 to 16 steps.} For the three CFM fields, SW2 decreases in all sets of both sizes. For the six fitted fields (one-grid fit and two-grid
fit), the mean change is negative at both sizes. With
1,024 samples, the standard deviation of the change across sets is
\VarSmallChangeSd{} on average over the six fields, and the mean change
is $\VarSmallChangeMean$. SW2 increases in \VarSmallFractionMin{}\% to
\VarSmallFractionMax{}\% of the sets, as a function of the field. At
least one of the six fields shows an increase in \VarSmallAnySix{}\% of
the sets, and all six fields show an increase together in
\VarSmallAllSix{}\% of the sets. The separate 1,024-sample set also shows all six increases. With 16,384 samples, the standard deviation of the
change is \VarMediumChangeSd{}, and SW2 increases in at most
\VarMediumFractionMax{}\% of the sets. On the set of 262,144
samples, SW2 decreases for all six fields
(Table~\ref{tab:large-bank}).

\emph{Change from 8 to 64 steps.} This change differs between the
fields. For the one-grid field with seed 17, the mean change is positive at both sizes. SW2 increases in \VarSmallGridLongFraction{}\% of the
sets of 1,024 samples and in \VarMediumGridLongFraction{}\% of the sets
of 16,384 samples. It also increases on the set of 262,144
samples. This agrees with the finite-witness intervals of
Section~\ref{app:large-bank-timing} for the same field. With 16,384
samples, the values for the other two one-grid fields are
\VarMediumLongOneB{}\% and \VarMediumLongOneC{}\%, and the values for
the three two-grid fields are \VarMediumLongTwoA{}\%,
\VarMediumLongTwoB{}\%, and \VarMediumLongTwoC{}\%.

With these fields and directions, increases from 8 to 16 steps occur
less often on sets of 16,384 samples than on sets of 1,024. The same
samples are used across fields within a set, so changes for the six
fields are correlated. The mean SW2 values in
Table~\ref{tab:set-variability} are also smaller on the larger sets.

\begin{table*}[ht]
\centering
\caption{Empirical SW2 of the nine two-dimensional fields on independent
evaluation sets. Entries are means across sets, with standard deviations
across sets in parentheses. The two increase columns give the percentages of sets in which SW2
increases from 8 to 16 and from 8 to 64 steps.}
\label{tab:set-variability}
\footnotesize
\begin{tabular}{@{}llrccccccc@{}}
\toprule
Training & Seed & Set size & SW2$_8$ & SW2$_{16}$ & SW2$_{64}$ & Change 8 to 16 & Increase & Change 8 to 64 & Increase\\
\midrule
CFM & 17 & 1,024 & 0.305 (0.029) & 0.243 (0.034) & 0.213 (0.040) & $-0.062$ (0.011) & 0\% & $-0.092$ (0.022) & 0\%\\
CFM & 17 & 16,384 & 0.280 (0.006) & 0.208 (0.006) & 0.169 (0.008) & $-0.072$ (0.003) & 0\% & $-0.111$ (0.007) & 0\%\\
\addlinespace
CFM & 29 & 1,024 & 0.266 (0.042) & 0.212 (0.049) & 0.193 (0.053) & $-0.054$ (0.014) & 0\% & $-0.073$ (0.026) & 0\%\\
CFM & 29 & 16,384 & 0.241 (0.012) & 0.177 (0.015) & 0.152 (0.018) & $-0.063$ (0.005) & 0\% & $-0.089$ (0.011) & 0\%\\
\addlinespace
CFM & 43 & 1,024 & 0.315 (0.042) & 0.251 (0.050) & 0.218 (0.056) & $-0.064$ (0.012) & 0\% & $-0.097$ (0.023) & 0\%\\
CFM & 43 & 16,384 & 0.297 (0.010) & 0.228 (0.012) & 0.189 (0.015) & $-0.070$ (0.004) & 0\% & $-0.108$ (0.008) & 0\%\\
\addlinespace
One-grid fit & 17 & 1,024 & 0.206 (0.050) & 0.201 (0.055) & 0.221 (0.056) & $-0.005$ (0.014) & 38\% & $+0.016$ (0.023) & 74\%\\
One-grid fit & 17 & 16,384 & 0.157 (0.014) & 0.147 (0.015) & 0.171 (0.014) & $-0.010$ (0.005) & 7\% & $+0.014$ (0.008) & 100\%\\
\addlinespace
One-grid fit & 29 & 1,024 & 0.260 (0.069) & 0.248 (0.073) & 0.259 (0.071) & $-0.012$ (0.012) & 18\% & $-0.001$ (0.020) & 48\%\\
One-grid fit & 29 & 16,384 & 0.244 (0.018) & 0.231 (0.020) & 0.241 (0.020) & $-0.013$ (0.004) & 0\% & $-0.003$ (0.007) & 30\%\\
\addlinespace
One-grid fit & 43 & 1,024 & 0.253 (0.054) & 0.246 (0.063) & 0.261 (0.067) & $-0.007$ (0.014) & 31\% & $+0.008$ (0.022) & 65\%\\
One-grid fit & 43 & 16,384 & 0.213 (0.011) & 0.202 (0.015) & 0.218 (0.016) & $-0.011$ (0.004) & 0\% & $+0.006$ (0.007) & 77\%\\
\addlinespace
Two-grid fit & 17 & 1,024 & 0.212 (0.048) & 0.199 (0.054) & 0.214 (0.056) & $-0.013$ (0.015) & 20\% & $+0.002$ (0.024) & 58\%\\
Two-grid fit & 17 & 16,384 & 0.166 (0.013) & 0.146 (0.015) & 0.163 (0.015) & $-0.020$ (0.005) & 0\% & $-0.003$ (0.008) & 33\%\\
\addlinespace
Two-grid fit & 29 & 1,024 & 0.265 (0.067) & 0.247 (0.072) & 0.253 (0.071) & $-0.018$ (0.012) & 6\% & $-0.012$ (0.020) & 29\%\\
Two-grid fit & 29 & 16,384 & 0.249 (0.018) & 0.230 (0.020) & 0.235 (0.020) & $-0.019$ (0.004) & 0\% & $-0.014$ (0.007) & 7\%\\
\addlinespace
Two-grid fit & 43 & 1,024 & 0.262 (0.053) & 0.248 (0.063) & 0.259 (0.068) & $-0.013$ (0.014) & 20\% & $-0.003$ (0.023) & 50\%\\
Two-grid fit & 43 & 16,384 & 0.223 (0.011) & 0.206 (0.014) & 0.217 (0.016) & $-0.017$ (0.005) & 0\% & $-0.006$ (0.007) & 20\%\\
\bottomrule
\end{tabular}

\end{table*}
